\section{Moments about the origin and the mean}
\textbf{Read:} Freund Ch4 \S3, pp.121--125; Wasserman \S\S3.1,3.3; Feller ChIX \S\S2--4. Practice set C.
\subsection{Raw moments}
The $r$th moment about the origin is $\mu'_r=\E(X^r)$, whenever finite. Thus $\mu'_0=1$, $\mu'_1=\mu$, and $\mu'_2=\E X^2$. Use LOTUS:
\[
\mu'_r=\sum_xx^rp_X(x)\quad\text{or}\quad\int x^rf_X(x)\dd x.
\]
Odd powers can cancel across a symmetric distribution; even powers are nonnegative. For a symmetric density about zero with finite relevant moments, odd moments vanish. A finite fourth absolute moment implies finite moments of orders one through three.
\begin{example}[A general moment formula]
For $f_X(x)=3x^2$ on $(0,1)$,
\[
\E(X^r)=\int_0^13x^{r+2}\dd x=\frac3{r+3}\quad(r>-3).
\]
Hence $\E X=3/4$, $\E X^2=3/5$, and $\E X^4=3/7$. One general formula replaces several nearly identical integrals.
\end{example}
\checkpoint{For $X$ uniform on $(-2,2)$, find $\mu'_1,\mu'_2,\mu'_4$.}{$0,4/3,16/5$, by symmetry and $\frac14\int_{-2}^2x^r\dd x$.}
\subsection{Central moments and variance}
The $r$th moment about the mean is $\mu_r=\E[(X-\mu)^r]$. Special cases are
\[
\mu_1=0,\qquad\mu_2=\sigma^2=\Var(X),\qquad\sigma=\sqrt{\Var(X)}.
\]
Expanding the square gives
\[
\Var(X)=\E(X^2-2\mu X+\mu^2)=\E X^2-\mu^2.
\]
Variance has squared units; standard deviation has the units of $X$. Also $\Var(X)\ge0$, with equality precisely when $X=\mu$ almost surely. Expanding higher powers gives
\begin{align*}
\mu_3&=\mu'_3-3\mu\mu'_2+2\mu^3,\\
\mu_4&=\mu'_4-4\mu\mu'_3+6\mu^2\mu'_2-3\mu^4.
\end{align*}
These formulas are optional conveniences; the definition always works.
\begin{example}[Variance of a square needs a fourth moment]
For the preceding $3x^2$ density,
\[
\Var(X)=\frac35-\left(\frac34\right)^2=\frac3{80}.
\]
For $Y=X^2$,
\[
\Var(Y)=\E X^4-(\E X^2)^2=\frac37-\left(\frac35\right)^2=\frac{12}{175}.
\]
Squaring $\Var(X)$ does not give $\Var(X^2)$.
\end{example}
\checkpoint{If $\E X=2$, $\E X^2=7$, $\E X^3=28$, find $\Var X$ and $\mu_3$.}{$3$ and $28-3(2)(7)+2(8)=2$.}
\subsection{Shifts, scales, and standardization}
For $Z=aX+b$, $Z-\E Z=a(X-\E X)$, so
\[
\E Z=a\mu+b,\qquad\Var Z=a^2\sigma^2,\qquad\operatorname{SD}(Z)=|a|\sigma.
\]
When $\sigma>0$, $Z=(X-\mu)/\sigma$ has mean zero and variance one; standardization does not generally make a distribution normal. The $r$th central moment of $aX+b$ is $a^r\mu_r$.
\checkpoint{If $\E X=10$, $\Var X=9$, $Z=5-2X$, find its mean, variance, and standard deviation.}{$-15$, $36$, $6$.}
